\subsection{Die zurückgewonnene Darstellung ist der ursprüngliche Träger}
Die sechzehn Gewichte zerfallen entlang der ursprünglichen ersten fünf und letzten drei Gitterrichtungen als
\begin{equation}
16\downarrow_{\mathrm{Spin}(10)\times SU(4)}=(10,1)\oplus(1,6)=V_{16}.
\end{equation}
Jede der $112$ ursprünglichen $D_8$-Wurzeln verbindet tatsächlich zwei Gewichtspaare dieser Darstellung. Auf P2 lauten ihre Ladungen
\begin{equation}
\{-\tfrac13\ (3\times),+\tfrac13\ (3\times),
-\tfrac12\ (2\times),+\tfrac12\ (2\times),0\ (6\times)\},
\qquad \tr_{V_{16}}Y^2=\frac53.
\end{equation}
Das rekonstruiert die vorher geometrisch gewonnene Darstellung mitsamt ihrer Wirkung. Es identifiziert diese Richtungen weiterhin nicht mit sechzehn vierdimensionalen Teilchensorten.

Auch die Überlagerung stimmt: Für $h=(z,-iI_4)$ ist $h^2$ der Kern der Einbettung nach Spin(16), während $h$ auf $V_{16}$ als $-I$ und auf dem gewählten Halbspinor als $+I$ wirkt. Somit ist die schon konstruierte Gruppe $K_f$ gerade die Paritätsüberlagerung der bosonischen Markierungsgruppe $K$. Eine allgemeine interne E8-Symmetrie wirkt innerhalb dieser Fermionisierung nur, wenn sie das gewählte $D_8$ erhält. Selbst dann können die Gruppenlifts die zentrale Fermionparität benötigen. Das ursprüngliche Gruppenwörterbuch bleibt maßgeblich.

\textbf{Die originale Viererstruktur lässt sich dabei ausdrücklich zurückgewinnen.} Für Cartanoperatoren $H_{0,i}$ mit Gewichten $\pm e_i$ ist der oben definierte duale Glue-Charakter
\begin{equation}\label{eq:inverseinternalglue}
G=\exp\!\left(i\pi\sum_{i=1}^5H_{0,i}\right).
\end{equation}
Auf den ersten zehn reellen Fermionen wirkt er mit $-1$, auf den letzten sechs mit $+1$. Im NS-Fockraum ist seine Wirkung deshalb $G_F=(-1)^{F_{10}}$ und $G_F^2=I$. Sein Spin-Lift besteht aber aus einer Halbdrehung in jeder der fünf ursprünglichen reellen Zweiebenen. Mit euklidischen Cliffordgeneratoren gilt
\begin{equation}
g=\prod_{j=1}^5(\gamma_{2j-1}\gamma_{2j}),\qquad g^2=(-1)^5=-1.
\end{equation}
Das zentrale $-1$ wirkt auf Vektoren trivial, auf Spinoren mit $-1$. Entsprechend tragen die ursprünglichen Spinorfelder $G$-Phasen $+i$ oder $-i$ und sehen Ordnung vier. Die $D_5$- und $D_3$-Ströme sehen $+1$, die gemischten $60$ Ströme $-1$. Genau daraus entsteht wieder die originale Klassenzählung $52+64+60+64$. Die Vorzeichenkonvention des Spin-Lifts kann $G$ und $G^{-1}$ vertauschen; die festgehaltene Orientierung bestimmt sie.

Damit ist $\mu_4$ eine tatsächliche Wirkung auf den ursprünglichen Sektoren. Der duale Glue-Charakter $G$, der gemeinsame Quotientsgenerator $h$ und die geometrische Vierteldrehung sind verschiedene, jetzt ausdrücklich beschriebene Operatoren. Ihre Unterscheidung erlaubt es, die alte Überlagerungsgeometrie an dieselbe Quelle anzuschließen, ohne ihre Wirkungen gleichzusetzen.

\subsection{Vakuum und Konformalzeit folgen zusammen}
Für die rekonstruierte Neveu--Schwarz-Darstellung sind die Moden und das Vakuum durch
\begin{equation}
\{\psi_r^a,\psi_s^b\}=\delta_{ab}\delta_{r+s,0},\qquad
r,s\in\Z+\tfrac12,\qquad \psi_r^a\Omega=0\quad(r>0)
\end{equation}
festgelegt. Der gemeinsame geerbte Stress-Tensor gibt
\begin{equation}\label{eq:inverseL0}
\boxed{L_0=\sum_{a=1}^{16}\sum_{r>0}r\,\psi_{-r}^a\psi_r^a,
\qquad L_0\Omega=0,\qquad c=8.}
\end{equation}
Die Realisierung des geraden CAR-Netzes, seiner Sektoren und der Virasorowirkung ist in~\cite{bockenhauercar} explizit konstruiert. Hier wird sie an die ursprüngliche TFPT-Markierung angeschlossen.

In Kreisradius-eins-Einheiten gilt $B_{\rm NS}=-i\partial_\theta$ mit antiperiodischer Randbedingung. Der positive Anregungsraum hat Energien $r=1/2,3/2,\ldots$ und $L_0=d\Gamma(h)$ mit $h=B_{\rm NS}|_{\mathcal K_+}$. Eine Zweitquantisierung von $|B_{\rm NS}|$ auf dem ganzen verdoppelten selbstdualen Raum würde die Moden doppelt zählen. Bei der Feldkonvention $B(e_r)=\psi_r$ ist die CAR-Zustandskovarianz die negative Fourierprojektion; der positive Anregungsraum und diese Kovarianz sind entsprechend konjugiert.

Die Antworten werden dadurch explizit. Auf der Ebene ist
\begin{equation}
\langle\psi^a(z)\psi^b(w)\rangle=\frac{\delta_{ab}}{z-w}.
\end{equation}
Für den positiven Halbzylinder, $\tau>0$, $\vartheta=\theta-\phi$ und Kreismaß $d\theta/(2\pi)$ liefert dieselbe positive Modensumme
\begin{equation}\label{eq:inversekernel}
K_{ab}(\tau,\vartheta)
=\delta_{ab}\sum_{n\ge0}e^{-(n+1/2)(\tau-i\vartheta)}
=\frac{\delta_{ab}}{2\sinh((\tau-i\vartheta)/2)}.
\end{equation}
Gerade höhere Majorana-Antworten sind Pfaffianen dieses Zweipunktkerns; ungerade verschwinden. Die alten C/X-Spinorfelder behalten ihre bereits berechneten Gittervertex-Antworten. Ihr gemeinsamer gerader Teil hat dieselben Vakuumantworten in beiden Erweiterungen.

Eine geometrische Vierteldrehung hat nun den tatsächlichen Spin-Lift
\begin{equation}
\rho_F=e^{i\pi L_0/2},\qquad \rho_F^4=(-1)^F,\qquad \rho_F^8=I.
\end{equation}
In der bosonischen Erweiterung $0\oplus s$ sind die Gewichte ganzzahlig, also $\rho_{E_8}^4=I$. Damit haben Double Cover, Vierteldrehung und Spinstatistik einen gemeinsamen Operatoranschluss. Geometrische Rotation, interner Glue-Charakter und eine Wilsonholonomie am rohen Nahtoperator sind trotzdem verschiedene Wirkungen. Eine originale Viertel-Wilsonholonomie muss als konkreter verdrehter Sektor identifiziert werden; sie wird nicht still zur NS-Randbedingung. Der festgehaltene konforme Stress-Tensor lässt innerhalb dieser Quelle keine unabhängig gewählte alternative Konformalzeit zu. Eine physische Zeiteinheit ist damit noch nicht kalibriert.

\subsection{Der lokale Rekonstruktionssatz: Was am Ursprung tatsächlich genügt}\label{sec:originkernelcriterion}
Der Rückweg liefert ein scharfes hinreichendes Ursprungskriterium. Angenommen, der ursprüngliche Prozess erzeugt eine unitäre, positiv graduierte chirale Vertexsuperalgebra mit eindeutigem Vakuumraum $V_0=\C\Omega$ und sechzehn reellen, gegenseitig graduiert lokalen Primärfeldern $\psi^a$ vom Gewicht $1/2$. Ihr positives Zweipunktprodukt sei nicht ausgeartet und in einer orthonormalen Basis normiert; ihre markierte Darstellung sei das originale $V_{16}$. Schließlich erzeugen diese Felder die betrachtete fermionische Quelle vollständig. Die zugehörige ursprüngliche Spinormarkierung bleibt erhalten.

\textbf{Dann sind die lokale Fermionenquelle, ihr Vakuum und ihr konformer Generator gemeinsam festgelegt.} Der Beweis benötigt keine angepasste Folge höherer Antworten:
\begin{enumerate}
\item Das singuläre Vertexprodukt $\psi^a_{(n)}\psi^b$ besitzt das Gewicht $1/2+1/2-n-1=-n$. Positive Energie verbietet $n\ge1$. Für $n=0$ ist wegen $V_0=\C\Omega$ nur das normierte Vakuumprodukt möglich. Deshalb ist der ganze singuläre Anteil zwingend
\begin{equation}
\psi^a(z)\psi^b(w)\sim\delta_{ab}(z-w)^{-1}.
\end{equation}
Gewicht und gewöhnliche graduierte Lokalität schließen hier weitere singuläre Felder aus.
\item Die Moden erfüllen somit CAR. Positive Energie erzwingt $\psi_r\Omega=0$ für $r>0$. Normalordnung bestimmt jedes höhere Vakuumwort; vollständige Erzeugung macht die CAR-Fockwörter dicht. Der Zustand ist nicht nachträglich frei wählbar.
\item Der freie Stress $T_F=-\frac12\sum_a:\psi^a\partial\psi^a:$ besitzt dieselbe Primärwirkung. Die Differenz seiner Moden zu einem zweiten solchen Stress kommutiert mit sämtlichen CAR-Moden. Sie bildet das Vakuum daher auf ein skalares Vakuum ab; Graduierung und Vakuumnormierung setzen diesen Skalar auf null. Durch Erzeugung verschwindet die Differenz auf der ganzen endlichen Fockdomäne. Es folgt genau Gleichung~\eqref{eq:inverseL0}. Ein unabhängiger zusätzlicher Sektor wäre mit vollständiger Erzeugung unvereinbar.
\item Der gerade Teil ist $(D_8)_1$. Die mitgeführte ursprüngliche Spinorerweiterung rekonstruiert $E_{8,1}$ mit den ursprünglichen Marken und demselben Stress-Tensor.
\end{enumerate}
Dies ist ein Rekonstruktionssatz in der angegebenen chiralen Feldklasse, keine Behauptung, P1/P2 hätten seine Voraussetzungen bereits geliefert. Seine entscheidende Leistung ist die Kompression: Ein \emph{lokaler markierter Feldkeim mit bewiesener Konformalwirkung} ersetzt die getrennte Suche nach kompletter Algebra, Zustand, höherer Korrelationshierarchie und Zeitgenerator. Eine Symbolklasse, ein bloßer Rang oder CAR zu einem einzelnen Zeitpunkt allein erfüllen dieses Kriterium noch nicht.

\subsection{Der mikroskopische Ausgangspunkt ist im Originalbestand bereits vorhanden}
Der Vertrag \file{microscopic-charged-car-limit/README.md} konstruiert am tatsächlich verwendeten QWZ-Zylinder einen lokalen geladenen CAR-Kanal mit beiden Adjungierten, seinem gefüllten Quellvakuum und kontrolliertem Zeitgrenzwert. Die tragende ursprüngliche Operatoridentität lautet für die obere Randzeile
\begin{equation}
h(p)r=-\sin(p)\,r+(1-\cos p)\,r_-,
\qquad r=\delta_{y,7}\chi_+,\quad r_-=\delta_{y,7}\chi_-.
\end{equation}
Sie liefert den chiralen linearen Anteil und einen quadratisch kleinen Mischterm. Daraus werden im Original alle festen endlichen Feldwörter samt Adjungierten, der positive Zustand und der Generatorgrenzwert auf einem gemeinsamen Energiekern gewonnen. Das ist ein tatsächlicher mikroskopischer Feldanschluss, auf dem weiter aufgebaut werden kann.

Die dortige ursprüngliche Holonomie wird ausdrücklich erhalten. Im Sektor $a=1/4$ ist der physische Quellgenerator
\begin{equation}
H_a=L_0+(a-\tfrac12)Q=L_0-\tfrac14 Q.
\end{equation}
Die im Vertrag \file{source-localized-spinor-20260919/PROOF.txt} bewiesene lokale Phasentransformation zum NS-Kernel ist nicht auf dem ganzen Kreis periodisch und entfernt diesen Term nicht. Der neue Rekonstruktionssatz bestimmt $L_0$; der vorhandene Quellgrenzwert zeigt konkret, wie die zusätzliche ursprüngliche Holonomie an die reale Zeit gekoppelt bleiben muss.

Die entscheidende noch zu liefernde Identifikation ist jetzt ebenso konkret: Der ursprüngliche markierte Nahtoperator muss den geometrisch bereits hergeleiteten $V_{16}$ tragen, einschließlich seiner gemischten $10\leftrightarrow6$-Wirkungen und der passenden Holonomie. Der vorhandene mikroskopische Vertrag liefert einen komplexen Randkanal. Sein nachfolgender Spinorvertrag setzt acht Kopien ausdrücklich voraus. Die transversale Breite acht ist keine Ableitung dieser acht Arten; der originale Kanaltest berechnet für alle transversalen Strombilineare einen führenden Gramrang eins. Dieser Befund richtet die Arbeit auf die tatsächliche Träger- und Operatorherleitung aus, statt denselben Ein-Kanal-Grenzwert erneut zu optimieren.

Auch die ursprüngliche Symmetrie markiert genau die erforderliche Verbindung: Unter $K_f$ sind $(10,1)$ und $(1,6)$ verschiedene irreduzible Summanden; allein diese Wirkung erzwingt keinen gemeinsamen Zeitgenerator beider Teile. Die gemischten $60$ $D_8$-Ströme verbinden sie tatsächlich. Wird ihre Wirkung am ursprünglichen Operator nachgewiesen, ist eine nur intern wirkende Kommutante auf $V_{16}$ skalar. Das bindet beide Blöcke an denselben skalaren Raum-/Zeitkern; es beweist noch nicht dessen konkrete Form. Genau diese Operatorwirkung muss aus dem Originalprozess kommen. Sie ist im rekonstruierten Ziel bereits explizit vorhanden.

\textbf{Dafür genügt zusammen mit der vorhandenen Untergruppenwirkung bereits ein wirklicher Mischstrom.} Der invariante Feldkern hat zunächst die Form $K=K_{10}P_{10}+K_6P_6$. Für einen vakuum- und zeitverträglichen reellen Kreuzgenerator gilt
\begin{equation}
X=\begin{pmatrix}0&A\\-A^{\mathsf T}&0\end{pmatrix},\quad A\ne0,
\qquad [K,X]=0\ \Longrightarrow\ K_{10}=K_6.
\end{equation}
Die Untergruppenbahn eines solchen Generators spannt die irreduziblen $10\otimes6$ Kreuzrichtungen. Ihre Lieklammern schließen zur ganzen $\mathfrak{so}(16)$-Wirkung. In der ursprünglichen, phasenfesten C/X-Tabelle ist ein passender komplexer Wurzelstrom schon ausdrücklich
\begin{equation}\label{eq:rawmixwitness}
\boxed{[C_0,X_{0,1}]_{\mathfrak e_8}=E_{-e_1-e_6}.}
\end{equation}
Der Koeffizient ist in der tatsächlich verwendeten Kokzykluskonvention $+1$. Zusammen mit dem adjungierten Wurzelstrom liefert er reelle Kreuzgeneratoren. Der unmittelbare Herkunftstest lautet deshalb: Lassen sich die lokalen $V_{16}$-Felder und diese markierte Stromwirkung im Originalprozess so realisieren, dass ihre Wardantwort den gemeinsamen Kernel erhält? Das ist eine konkrete verbindende Operatorfrage; sie verlangt nicht die unabhängige Anpassung sämtlicher späterer Antworten. Die Identität~\eqref{eq:rawmixwitness} ist im gewählten Quellenmodell berechnet, ihr mikroskopischer Rohfeld-Intertwiner bleibt zu konstruieren.

\subsection{Wie das den Weg zur fundamentalen Lösung verändert}
Der Rücklauf
\begin{equation}
(E_{8,1};D_8,s)\longmapsto\mathcal F_{16}\ \text{mit Markierung }s
\longmapsto(E_{8,1};D_8,s)
\end{equation}
ist innerhalb der gewählten Quellenklasse geschlossen. Die schon vorhandene rekursive Rückbestätigung erhält damit ein Feld-, Vakuum- und Zeitgegenstück. Dieser Rücklauf liefert ein präzises Ziel für die Herkunftsherleitung; er ersetzt sie logisch nicht.

Am Originalprozess ist nun genau der in Abschnitt~\ref{sec:originkernelcriterion} bezeichnete Keim nachzuweisen: seine tatsächlichen lokalen Felder auf dem bereits konstruierten $V_{16}$, ihre halbe konforme Gewichtswirkung und ihre vollständige Erzeugung. Dazu gehört die originale Holonomie-/Spinorsektorabbildung. Diese Bedingungen sind gemeinsam zu bearbeiten. Neue Oszillatoren, zusätzliche Clockraten oder eine andere Gitterwelt würden diesen Nachweis nicht liefern.

Für den Universalraum bedeutet das: Sein positiver Geschichtskern ist nach dieser Identifikation bereits durch dieselbe Feldregel bestimmt. Auf dem gemeinsamen geraden Teil stimmen die beiden Darstellungen überein; geladene Geschichten müssen ihren NS-/Spinorsektor behalten. Zuses Komposition und ein kausaler Ereignisgraph können diese Felder und Geschichten organisieren, sofern sie ihre Produkte, Sektoren und Antworten erhalten. Sie müssen die Grundregel dann nicht erneut frei wählen. Die tatsächliche Clock-Auslese und die vierdimensionale lokale Realisierung werden weiter an dieser einen Quelle angeschlossen; eine chirale Rekonstruktion allein bestimmt sie noch nicht.

\picturetext{Bislang sah es so aus, als müssten wir für ein Instrument Saiten, Grundzustand, Klanggesetz und Takt einzeln finden. Der Rückweg zeigt: Beim vorhandenen Instrument gehören diese Teile bereits zwingend zusammen. Wir kennen jetzt die kurze lokale Regel, aus der sie gemeinsam entstehen. Die fundamentale Aufgabe ist, genau diese Regel am ursprünglichen Nahtprozess nachzuweisen. Danach müssen die räumlichen Messungen als wirkliche Betätigungen desselben Instruments erkennbar werden.}

\subsection{Ausführbare Prüfung und ihr genauer Umfang}
Die bestehende Quellenimplementierung wurde um \file{reconstruct_marked_fermion_source} ergänzt. Sie benutzt die originalen E8-Wurzeln und Feldmarkierungen, berechnet den geraden Glue-Fixpunkt, die $10+6$-Wirkung und die P2-Ladungen. Die Tests prüfen die Glue-Graduierung an sämtlichen ursprünglichen Wurzelklammern und die tatsächliche Wirkung der $D_8$-Wurzeln auf den sechzehn Gewichten.

Zwei unabhängig berechnete Charaktere stimmen überein: die NS-Fermionenbesetzung $\prod_{n\ge0}(1+t^{2n+1})^{16}$ und der ungerade Gittercharakter mit acht bosonischen Oszillatoren. Für $t^{2L_0}$ beginnen sie mit $1+16t+120t^2+576t^3+2076t^4$. Der mit originaler Chiraliät berechnete Spinorsektor beginnt mit $128t^2+2048t^4$. Gerade NS-Quelle plus Spinor ergibt $1+248t^2+4124t^4+\cdots$ und wurde bis $L_0=6$ zusätzlich gegen $E_4/\prod_{n\ge1}(1-q^n)^8$ geprüft. Diese endlichen Vergleiche kontrollieren Implementierung und Markierung; der allgemeine Rekonstruktionssatz steht im Beweis darüber. Er ist hier nicht als neuer Lean-Beweis ausgegeben.

\section{Wie die Gesamtlösung konkret erreichbar würde}\label{sec:completion}

\key{Die Aufgabe bleibt die gemeinsame \textbf{Herkunft, Prozessrekursion und physische Realisierung}. Die vorhandenen physikalischen Herleitungen aus Abschnitt~\ref{sec:physicalconstraints} geben dabei konkrete Auswahlbedingungen vor. Vorrang hat ihre simultane Realisierung am ursprünglichen markierten Flavor- und Nahtoperator mit derselben Determinantenfamilie. Der neue E8-Kandidat, sein Clockinstrument und seine Raumzeitfortsetzung müssen diese Antworten gemeinsam erhalten. Eine weitere Rechnung nur im symmetrischen Quellmodell genügt dafür nicht.}

\subsection{Erste Brücke: Die Quantisierung aus dem Ursprung auswählen}
Die Zielaussage lautet: Die ursprünglichen P1/P2-Daten samt den tatsächlich schon begründeten Markierungen erzwingen oder charakterisieren die minimale positive lokale Level-1-Vakuumquantisierung des markierten Gitters innerhalb einer ausdrücklich definierten Vergleichsklasse.

Der Auswahlsatz für die unitäre affine Vergleichsklasse ist in Abschnitt~\ref{sec:levelselection} jetzt konkret geführt: Eine unabhängig begründete Zentralzahl höchstens acht oder ein geeigneter langer Wurzelstrom mit Doppelmoden-Nullrelation erzwingt Level eins. Der nächste Herkunftsnachweis betrifft deshalb die ursprüngliche Realisierung genau dieses Stroms beziehungsweise dieser Zentralzahlgrenze. Primitive U1-Windung, Rang acht und die normierte Fünferrekursion dürfen nicht einzeln als bereits hinreichende Ersatzbegründung verwendet werden.

Abschnitt~\ref{sec:onespinorcurrent} präzisiert den globalen Architekturtest: Ein einziger Spinorstrom vom Gewicht eins genügt; eine vorherige E8-Auswahl oder vollständige Stromerzeugung ist unnötig. Abschnitt~\ref{sec:doubletwistbridge} identifiziert im bedingten Achtkanalrahmen seinen führenden Twist-Fusionskoeffizienten. Der vorhandene Halb-Twist-Vertrag liefert bereits den mikroskopischen Randbedingungsschatten und den benötigten Ladungs-Carry. Nachzuweisen ist ihre gemeinsame lokale Realisierung auf dem ursprünglichen $V_{16}$-Träger, einschließlich Vakuum, Adjungiertem und Energiegrenzen. Diese Forderung wird nicht durch zusätzliche deklarierte QWZ-Kopien erfüllt.

Abschnitt~\ref{sec:microscopicpairbridge} schließt dazu erstmals den distributionswertigen Gewicht-1-Paarkern der deklarierten Achtkopienquelle analytisch an die volle mikroskopische Antwort an. Die noch zu erbringende Ursprungsaussage ist damit nicht erneut eine abstrakte Existenz irgendeines Kontinuums: Sie betrifft Gleichung~\eqref{eq:originalmultiplicitybridge} und die lokale geladene Feldwirkung samt ursprünglicher Deckmarkierung. Die vorhandenen Alpha-, Flavor- und Massenauslesungen sind anschließend an genau diese transportierte Verbindung und ihren Zustand zu binden.

Die rekursiven P1/P2-Rückbestätigungen gehören in diesen Satz. Sie dürfen weder ignoriert noch so verwendet werden, als seien sämtliche Voraussetzungen ihres Abschlussrahmens bereits hergeleitet. Eine bloße Suche nach beliebigen alternativen Kanälen ohne die ursprünglichen Marken wäre hier kein Fortschritt.

Die Fortsetzung in Abschnitt~\ref{sec:nativeanomaly} liefert dafür einen direkteren Eingang: Der ursprüngliche globale $V_{16}$-Träger ist konstruiert. Wird seine tatsächliche gleichchirale Dirac-/Pfaffian-Realisierung nachgewiesen, folgt $c=8$ aus dem Familienindex, ohne zuvor E8 in die Rechnung einzusetzen. Der echte CAR-/Projektionsanschluss bestimmt dann den vollständigen reinen Zustand; die markierte lokale Spinorerweiterung schließt an E8 an. Die bordismustheoretische Route aus Abschnitt~\ref{sec:bordism16} prüft ergänzend den Abstieg der Antwort ohne Spinwahl. Der K3-Diracindex oder ein eingesetzter Komponentenfaktor ersetzt keinen dieser tatsächlichen Operatoranschlüsse.

\subsection{Zweite Brücke: Ein vollständiger Prozess unter Rekursion}
Die Zielgleichungen sind
\begin{equation}
\widetilde A V=VA,\qquad \widetilde U(t)V=VU(t),
\end{equation}
für die gemeinsam zugelassenen geladenen Operationen und die tatsächlich identifizierte Zeit. Die vorhandenen $W$-Intertwiner, Hecke-Sektorprojektoren und geladenen Carry-Verbraucher sind konkrete Bausteine. Der Nachfahrenturm und die bei endlichen Abständen sichtbaren Restsektoren dürfen dabei nicht abgeschnitten werden, wenn die Behauptung vollständige Antworttreue lautet.

Für radiale Wurzelfeldgeschichten ist diese Verbindung jetzt ausdrücklich konstruiert: $Q,A,P(t)$ bestimmen den vollen Vektor, seine zulässige Felderweiterung und den positiven Kernel; die Möbiusformel erhält den vollständigen Nachfahrenturm. Eine vorgegebene historische Vergröberungsregel muss auf diesen Daten als antworterhaltende Abbildung identifiziert werden. Der nächste Schluss ist somit kein beliebiges Anbauen von Restregistern, sondern der Nachweis, dass genau die vorhandene Rekodierung denselben Zustand und dieselben Operationen erhält. Das ist noch kein allgemeiner kausal normierter Detektorprozess.

Die Replica-Konstruktion präzisiert nun die Clockseite: Beide Beträge $2/3$ und $1/3$ treten in der signierten Vierpunktantwort auf. Die positiven dreidimensionalen Populationen erfordern jedoch eine wirkliche Auslese gemäß Gleichung~\eqref{eq:nativeclockmap}. Die unzulässige Gleichsetzung dreier Wickterme mit drei Zuständen und die reine $L_0$-Zeitumkalibrierung sind ausgeschlossen. Damit ist klar, welche Art von Prozessabbildung noch sinnvoll zu suchen ist.

\subsection{Dritte Brücke: Dieselbe Quelle in der physischen Raumzeit}
Die Originalverträge benennen bereits den richtigen gemeinsamen Prüfgegenstand. In passender Notation lautet er
\begin{equation}\label{eq:physical}
\operatorname{Res}_\Sigma W_{4D}[j]=W_{E_{8,1}}[j],
\end{equation}
zusammen mit
\begin{equation}
\operatorname{Res}_\Sigma\det D_{4D}\cong\det D_{\rm seam}
\quad\text{als Linienbündel mit Verbindung}.
\end{equation}
Zu spezifizieren sind die Einsetzungs- und Restriktionsabbildung, die zulässigen Quellen $j$, die Feldnormalisierung, P2-Ladungen, Clockmarkierungen, Zustand, Zeit und analytischen Gebiete. Die Gleichheit eines ausgewählten Zahlenwertes oder bloßer Determinantenbeträge reicht nicht.

Die entsprechenden Originalkennungen sind
\file{TFPT4D.LATTICE.ACTION.01},
\file{SEAM.DETLINE.UNIFICATION.01} und
\file{FTRANSFER.GENERATING.01}.
Die vorhandene endliche $\Z_6$-Funktionalroute und die $2+1$-dimensionale QWZ/$U(1)$-Determinantenroute sind dafür Vorarbeiten. Sie enthalten noch nicht den vollständigen chiralen $3+1$-$E_8$-Anschluss.

\subsection{Der gemeinsame Prozess bestimmt noch nicht seine räumliche Zuordnung}
Der stärkere Originalvertrag \file{BULK.PRIME_COMPLETION.01} verlangt ausdrücklich eine durch die Randoperationen erzeugte Realisierung ohne unsichtbaren Zusatzsektor. Unter dieser Bedingung dürfen beliebige angehängte Zuschauer nicht als Gegenbeweis gegen Eindeutigkeit verwendet werden. Sind auf einer zyklisch erzeugenden Wortalgebra sämtliche Antworten
\begin{equation}
D(A,B)=\omega(A^*B),\qquad C_t(A,B)=\omega(A^*\alpha_t(B))
\end{equation}
gegeben, wobei $\alpha_t$ eine $*$-Automorphismengruppe mit $\omega\circ\alpha_t=\omega$ ist, bestimmt $D$ die GNS-Darstellung. Die Zustandsinvarianz macht die Wortverschiebung $[A]\mapsto[\alpha_t(A)]$ isometrisch; die Gruppeninversen machen sie unitär. $C_t$ bestimmt sämtliche Matrixelemente dieser Entwicklung auf der dichten Menge der Wortzustände. Unter starker Stetigkeit ist damit auch ihr Generator bis auf die entsprechende unitäre Identifikation festgelegt. Das ist die positive Aussage des ursprünglichen KMS-Kohärenzvertrags; innerhalb dieses vollständigen Datenvertrags bleibt kein unabhängig frei wählbarer Hamiltonoperator übrig.

Noch nicht enthalten ist dadurch die Abbildung
\begin{equation}
O\subset\R^{1,3}\longmapsto\mathcal A_{\rm bulk}(O).
\end{equation}
Sie ordnet den Operatoren ihre räumliche Zugänglichkeit zu und muss mit Symmetrie, Lokalität, Quelle und Wechselwirkung gemeinsam verträglich sein. Bei der vorhandenen holomorphen E8-Netzstruktur liefern zwei alternierende Intervalle und ihre Vierteldrehung aufgrund der Zweiintervall-Dualität lediglich $M$ und $M'$. Ihr Schnitt ist $\C I$; das Vakuum erzeugt daraus nur eine Linie. Dies ist keine standarde Vierkeilkonfiguration für eine vierdimensionale Modularrekonstruktion. Der passende Rekonstruktionssatz benötigt gerade eine bestimmte gemeinsame Stellung der Algebren als Voraussetzung~\cite{kahlerwiesbrock}. Vier Marken allein liefern sie nicht.

\subsection{Konforme Kreiszeit und physische Translation müssen getrennt werden}\label{sec:physicaltimeaudit}
Auch die Formulierung, der vorhandenen Konformalzeit fehle lediglich eine physische Einheit, wäre für den exakten Poincar\'e-Vertrag zu schwach. Auf dem E8-Vakuumraum gilt $e^{-2\pi iL_0}=I$. Jede affine Wahl $H=aL_0+bI$ mit $a\ne0$ besitzt daher eine skalare Periode. Beim fermionischen Cover verdoppelt sich die Periode; das Argument bleibt bestehen.

\textbf{Satz.} In einer global Poincar\'e-kovarianten unitären Darstellung mit Translationen $V(x)$ erzwingt eine skalare Periode $V(Te_0)=\zeta I$, $T>0$, die Trivialität sämtlicher Translationen.

\textbf{Beweis.} Gegenläufige Boosts entlang der $i$-ten Raumrichtung liefern
\begin{equation}
V(T\cosh\eta,\,\pm T\sinh\eta\,e_i)=\zeta I.
\end{equation}
Ihr Quotient ergibt $V(0,2T\sinh\eta\,e_i)=I$. Weil $\sinh\eta$ alle reellen Werte annimmt, sind sämtliche Raumtranslationen trivial. Der Quotient einer geboosteten Periode mit der ursprünglichen ergibt danach $V(T(\cosh\eta-1),0)=I$. Diese Zeitabstände durchlaufen $[0,\infty)$; mit den Inversen sind auch alle Zeittranslationen trivial. Ein Vakuumshift $b$ oder eine Kalibrierung $a$ verhindert den Schluss nicht.

Dieser Satz schließt die direkte Identifikation der kompakten Quellenrotation mit einer nichttrivialen Minkowski-Zeittranslation aus. Er widerlegt keine approximative Raumzeit in einem Grenzübergang und keine gekrümmte kosmologische Realisierung ohne globale Poincar\'e-Symmetrie. Er gilt für die exakte Identifikation im genannten Originalvertrag.

Die vorhandene Möbius-Darstellung enthält bereits einen möglichen anderen Zeittyp. In einer Cayley-Konvention ist
\begin{equation}
P_{\rm line}=L_0-\tfrac12(L_1+L_{-1})
\end{equation}
der positive Generator der parabolischen Translation auf dem aufgeschnittenen Kreis. Er stammt aus demselben Stress-Tensor und derselben Quelle, ist aber keine affine Umrechnung von $L_0$. Die Wahl des ausgezeichneten Unendlichkeitspunktes gehört zu dieser Identifikation. Die Unterscheidung von Rotations-, Translations- und Sonderkonformgenerator und die zugehörige Positivität sind Standard in der Möbius-Darstellungstheorie~\cite{koesterconformal,morinellirehren}. Damit ist eine nichtperiodische chirale Bewegung vorhanden; ihre Identifikation mit physischer Zeit und die fehlenden transversalen Raumwirkungen folgen daraus noch nicht. Der vorhandene Vertrag \file{source-local-line-20260920/LOCAL_LIMIT_REVIEW.txt} liefert außerdem einen kontrollierten lokalen Großradiusgrenzwert. Dieser verändert die Spektralstruktur und ist durch den Satz über eine feste periodische Darstellung ausdrücklich nicht ausgeschlossen.

Die in derselben Primärarbeit beschriebenen räumlichen Deformationen zeigen, dass ein fester Zeitprozess durchaus verschiedene räumliche Realisierungen tragen kann. Dort werden jedoch eine passende vierdimensionale Einteilchendarstellung und ihre räumlichen Generatoren verwendet; die Konstruktion betrifft freie Felder. Sie darf weder als Beweis eines allgemeinen Dimensionshindernisses noch als bereits vorhandene wechselwirkende TFPT-Raumzeit ausgegeben werden.

\subsection{Warum die 24 Felder die Aufgabe wesentlich verkleinern}
Die 24 ursprünglichen Felder zusammen mit ihren Adjungierten erzeugen die gesamte Stromalgebra. Wenn eine physische Realisierung ihre vollständigen Level-1-Relationen, ein gemeinsames positives zyklisches Vakuum, analytische Lokalität und lokale Erzeugung nachweist, ist der erzeugte Vakuumsektor dadurch stark festgelegt. Unter den passenden Rekonstruktionsvoraussetzungen bestimmen diese Daten dieselben höheren Randantworten und dasselbe Netz.

Das ist die entscheidende Kompression des Problems: Statt tausend Antworten unabhängig anzupassen, werden die Erzeuger, ihr vollständiges Produktgesetz, der Zustand und die Zeit identifiziert. Dieser Satz allein bestimmt noch keine eindeutige vierdimensionale Erweiterung. Er macht aber ihren Randanschluss zugleich umfassend und überprüfbar.

Die Sechserantwort $1/576$ ist dabei ein notwendiger strenger Kontrollfall für jede behauptete Realisierung \emph{dieser} Quelle. Sie ist weder allein hinreichend noch ein Widerlegungstest für Modelle, die diesen Quellenanschluss gar nicht behaupten. Ihr Nutzen besteht darin, eine bloße Paarungs- oder Registerimitation unmittelbar zu erkennen.

\subsection{Welche physikalischen Zweige danach am selben Objekt hängen}
\begin{longtable}{L{0.21\linewidth}L{0.38\linewidth}L{0.33\linewidth}}
\toprule
Zweig & Gemeinsames Objekt & Entscheidende noch benötigte Identifikation\\\midrule\endhead
Eichstruktur & Geladene Ströme, Wardidentitäten, lokale Gauss-Wirkung & Interne Symmetrie als tatsächliche lokale Eichwirkung\\
Chirale Materie & Physische Feldoperatoren und ihre Spektren & Chirale Raumzeitwirkung und Spiegelentkopplung\\
Flavor/Massen & Pole und Mischungsantworten derselben Mehrpunktfunktionen & Yukawa-Dynamik und Maßstab, nicht nur Rangzählung\\
$\alpha$ & Elektromagnetische Wardantwort und Normierung & Herkunft genau der registrierten Selbstkonsistenzgleichung\\
Gravitation & Metrische Variation desselben physischen Funktionals & Dynamische geometrische Antwort und Quantisierungsbereich\\
Kosmologie & Mehrzeitfunktional einschließlich Anfangszustand & Physische Zeit, CP-/Nichtgleichgewichtsantworten und Transfer\\\bottomrule
\end{longtable}

Für Nichtgleichgewichtsfragen gehört der Anfangszustand ausdrücklich zum Schwinger--Keldysh-Objekt. Ein konformer Vakuumkorrelator allein ist kein kosmologischer Anfangszustand. Die bekannten bedingten Vorhersageformeln behalten ihren Wert als gemeinsame Zielantworten, ihre Transferannahmen müssen jedoch im selben physischen Objekt erklärt werden.

\subsection{Was nach diesen Prüfungen tatsächlich herzuleiten bleibt}
Die Suche setzt am jetzt bewiesenen Rekonstruktionskriterium aus Abschnitt~\ref{sec:originkernelcriterion} an. Die originale Naht muss die tatsächlichen lokalen $V_{16}$-Felder, ihre halbe konforme Gewichtswirkung und vollständige Erzeugung liefern. Daraus folgen CAR-Kernel, reiner Zustand, Konformalzeit und die markierte Erweiterungsroute gemeinsam. Der Rückweg von E8 liefert bereits das genaue Zielobjekt; die Herkunftsherleitung muss es unabhängig am ursprünglichen Prozess wiederfinden. Der skalare DtN-Wert $|k|$ aus den älteren Verifikationen bestimmt dabei noch nicht das Cliffordmodul oder diese Feldwirkung.

Auf der Clockseite wird die tatsächliche Cusp-Auslese in derselben Quelle gesucht, nicht ein weiterer Kanal mit vorab eingesetzten Eigenwerten. Ihr Erfolgskriterium ist die Operatorgleichung~\eqref{eq:nativeclockmap} einschließlich der ursprünglichen sechs Kompositionen. Die schon bekannte freie $S_3$-Kanalfamilie darf diese Herkunftsfrage nicht durch erneute Parameterwahl verdecken. Die neue gemeinsame Klassifikation schließt außerdem $q=1/3$ als autonome zeitunabhängige Markoventwicklung derselben Register aus. Deshalb muss ein vorgeschlagener Quellenanschluss seine Records und Zwischenoperationen mitführen. Ein bloßer Ein-Schritt-Abgleich $\mathcal E\circ\alpha_\delta\circ\iota=\Phi_q$ reicht nicht: Die identischen Abbildungen müssen auch die tatsächlich beanspruchten geordneten Mehrzeitantworten erhalten. Der lokale Puls liefert hierfür eine bereits definierte Quellenintervention statt einer zusätzlichen frei gewählten Dynamik.

Für die Raumfrage reicht die kanonische Vierteldrehung der schon vorhandenen Netze nicht: Die vorstehende Prüfung liefert keine geeignete vierdimensionale Stellung. Gesucht ist eine aus den Ursprungsoperationen begründete zusätzliche relative Lokalisierung mit tatsächlich transversalen Wirkungen. Sie muss die geordneten Antworten erhalten und die Raumzeitsymmetrie zusammen mit dem richtigen Translationstyp tragen. Neue räumliche Argumente dürfen nicht stillschweigend in eine vorausgesetzte $4D$-Masterwirkung eingesetzt werden.

Die so verbundenen Feld-, Zustands-, Ladungs- und Zeittests müssen gemeinsam bestehen. Die neue Realisierung muss dabei die bereits vorliegenden Massen-, Yukawa-, $\alpha$- und Newtonantworten aus Abschnitt~\ref{sec:physicalconstraints} unter ihren jeweiligen Originalvoraussetzungen reproduzieren. Die Fortsetzung hat mehrere vorher offene strukturelle Anschlüsse konstruiert, jedoch diese verbleibenden Operatoridentitäten und die vollständige Raumzeitphysik noch nicht bewiesen.

\section{Neue Anwendungen: Wo die Verbindung bereits produktiv wird}

``Neu'' hat hier zwei verschiedene Bedeutungen. Eine Anwendung kann für TFPT neu sein, obwohl das mathematische Verfahren bekannt ist. Wissenschaftliche Erstmaligkeit oder ein bahnbrechender Leistungssprung verlangt zusätzlich einen Vergleich mit der einschlägigen Literatur und den stärksten vorhandenen Verfahren. Die folgenden Anwendungen ergeben sich konkret aus den Sitzungsbefunden.

\begin{longtable}{L{0.18\linewidth}L{0.38\linewidth}L{0.36\linewidth}}
\toprule
Bereich & Konkrete Anwendung & Woran ihr Nutzen erkennbar wäre\\\midrule\endhead
Geometrie & Aus derselben markierten Quelle geladene Blockabbildungen und effektive Ausbreitung berechnen & Korrekte Antworten unter geänderter Zellgeometrie, einschließlich harmonischem Korrektor\\
Topologie & Holonomie und Deckphasen als endlichen beziehungsweise rationalen fortsetzbaren Zustand speichern & Gleiche nackte Gitterwege werden nur dann identifiziert, wenn alle erlaubten späteren Antworten gleich bleiben\\
Mathematik & Ein symbolischer Rechner für gemischte Stromwörter, verbundene Antworten und zulässige Fusionswege & Neue exakte Antworten ohne neue fitbare Ereignisregeln; automatische Gegenfälle gegen falsche Faktorisierung\\
Quanten\-information & Geladene Codes und Records mit nachgewiesenem Verbraucher bauen & Intertwiner erhalten die tatsächliche Operation; Leakage und Nachfahrensektoren werden quantitativ kontrolliert\\
Physik & Kandidaten für den gemeinsamen Rand-/Raumzeitanschluss an denselben höheren Antworten prüfen & Ein Kandidat besteht Ladungs-, Zeit-, Positivitäts- und Mehrpunktprüfung zusammen\\
Experiment\-design & Anhand von Ladung und Operatorstruktur gezielt höhere gemeinsame Antworten suchen & Weniger Versuche bis zu einem trennenden Befund als bei zufälliger oder rein paarweiser Auswahl\\
KI/Gedächtnis & Zustände nach späterer Handlungswirkung komprimieren und Gegenbeispiele als nötige Erinnerung behalten & Bessere wiederverwendbare Vorhersagen und Entscheidungen bei gleichem Speicher und Interaktionsbudget\\\bottomrule
\end{longtable}

Die besonders produktive mathematische Anwendung ist die automatische Entdeckung \emph{falscher Vereinfachungen}: Ein Rechner findet das kleinste gemischte Wort, an dem zwei vermeintlich gleiche Darstellungen unterschiedliche Antworten liefern. Das verbindet Beweisarbeit, Experimentwahl und Programmoptimierung unmittelbar.

Eine spekulative physikalische Anwendung ist der gezielte Entwurf von Systemen mit unterdrückten niedrigen und sichtbaren höheren Antwortkoeffizienten. Der Sechserzeuge liefert dafür ein exaktes algebraisches Muster. Die Umsetzung in ein konkretes Material oder Experiment würde jedoch eine zusätzliche Kopplungs-, Präparations- und Messabbildung erfordern. Das Dokument behauptet keine bereits konstruierte Hardware und keine neue fundamentale Sechskörperkraft.

\section{Welche Algorithmen daraus entstehen könnten -- und was daran bahnbrechend wäre}\label{sec:algorithms}

\key{Der stärkste gemeinsame Algorithmusgedanke ist ein \textbf{Compiler für zukünftiges Verhalten}: Er gewinnt aus tatsächlichen Antworten die nötigen Zustandsunterschiede, speichert sie in einer geeigneten Struktur und übersetzt zulässige Operationen in ein kleines, nachweislich antworttreues Programm.}

Dieser Gedanke verbindet die fünf folgenden Bausteine. Sie sollten als ein Arbeitsweg entwickelt werden, nicht als fünf neue Modellwelten. TFPT liefert dafür eine ungewöhnlich konkrete Quelle exakter Antworten und Gegenfälle. Die Universalraumidee liefert das Kriterium, welche Information erhalten bleiben muss. Symmetrie liefert eine Ordnung dieser Information; Verifikation sichert die erlaubten Umformungen.

\subsection{Der Neuheitsmaßstab und ein bereits vorhandener Gegenvergleich}
Zustände durch zukünftige Tests zu definieren ist bekannt. Aktive Automatenlerner, Predictive State Representations, Observable Operator Models und gewichtete Automaten verfolgen verwandte Ziele. Auch Verbindungen zu Tensor-Netzen sind publiziert \cite{littman,thon,adhikary,lsharp}. Die bloße Zusammenführung dieser Begriffe ist daher kein neuer Universalalgorithmus.

Ein historischer eigener Vergleich vom 25. September wurde für dieses Dokument am Originalbericht nachgelesen. Er untersuchte vier endliche Aufgabenfamilien, zehn Seeds, zwölf Aktionen und ein Budget von 30000 Operationen pro Lauf. Der reparierte TFPT-inspirierte Causal Compiler rekonstruierte 37 von 40 vollständigen Minimalmodellen; $L^\#$ mit derselben Heuristik erreichte 40 von 40. Die frühere Variante lag bei 30 von 40.

\begin{center}
\begin{tabular}{lrr}
\toprule
Familie & Reparierter Compiler & $L^{\#}$ mit Heuristik\\\midrule
Shutter & 20554.5* & 7675\\
Dense & 6159.5 & 3349.5\\
Stack & 11109 & 8375.5\\
SQLite & 12016 & 9327.5\\\bottomrule
\end{tabular}
\end{center}

Die Tabelle zeigt den Median der Operationen bis zum ersten vollständigen Minimalmodell. Eine Operation ist ein Reset oder eine Aktion einschließlich Readout. Beim mit Stern markierten Shutter-Wert werden ungelöste Fälle mit dem vollen Budget 30000 gezählt; dies ist ein gedeckelter Score.

Die reine Lernschleifenzeit über alle vierzig Läufe bis zum gemeinsamen Budgetende betrug dagegen 3,368 Sekunden für den reparierten Kern und 27,558 Sekunden für $L^{\#}$ mit Heuristik. Das ist keine Zeit bis zu einer selbst erkannten vollständigen Lösung. Weniger Umweltoperationen bedeuten also nicht automatisch weniger Rechenzeit.

Die Verbesserung des eigenen Lerners beruhte auf bekannter binärer Gegenbeispielanalyse. Bei einer passwortartigen Aufgabe ohne informative Zwischenantworten löste keiner der sieben Lerner die zehn Fälle im Budget. Ein SQLite-Beispiel zeigte zwar, dass gleiche sichtbare Tabellen nach einem späteren Rollback verschieden reagieren können; gewöhnlicher Zustands- beziehungsweise Transaktionsspeicher kann diese Information jedoch ebenfalls erhalten.

Der Originalvergleich umfasste 280 Läufe und unabhängige Modell-, Zeugen- und SQLite-Kontrollen. Er ist ein datierter früherer Befund, kein in dieser Sitzung erneut ausgeführter Benchmark und kein aktueller Hylæan-Test. Er verbietet, denselben Automatenlerner nun unter geometrischem Namen als Durchbruch zu vermarkten. Eine neue Methode muss einen anderen, konkreten Zusatz leisten: etwa positive operatorwertige Prozesse kontrolliert komprimieren, neue Zukunftsfehler mit weniger Eingriffen finden oder antworttreue Programme mit überprüfbaren Fehlergrenzen erzeugen.

\subsection{Baustein 1: Ein schneller Compiler für das vollständige Quellgesetz}
Die Forschungsfortsetzung hat hierfür bereits einen konkreten Teil umgesetzt: Für sämtliche endlichen C/X-Wurzelwörter wertet \file{evaluate_lattice_word} das ursprüngliche Gesetz durch quadratisch viele exakte Paaroperationen aus. Die tatsächlichen Kokzyklusphasen und geladenen Erweiterungssektoren bleiben erhalten. Die Ward-Rekursion ist der unabhängige Vergleich. Das ist eine im TFPT-Rechner nutzbare Anwendung einer bekannten Gitterformel; weder wissenschaftliche Erstmaligkeit noch ein gemessener allgemeiner Geschwindigkeitsrekord wird behauptet. Die Einschränkung auf Wurzelfeldeinsetzungen und die wachsende rationale Bitgröße gehören zum Leistungsumfang.

Die bereits ausgeführte Ward-Regel ist der Rechenkern. Ein weiterentwickelter Compiler würde gleiche Teilwörter nur einmal auswerten, ladungsunmögliche Beiträge vorab entfernen und tatsächlich äquivalente Anfragen kanonisieren. Für den bewiesenen harmonischen $X$-Sektor darf er den schnellen Weg über $W_5\otimes W_4$ verwenden. Für gemischte $C/X$-Wörter muss er auf die vollständige Regel zurückfallen.

Die konkrete Neuheitschance liegt in einer nachgewiesen guten Auswahl der Auswertungsreihenfolge, die Ladung, Sektor, Sparsität und Kosten gemeinsam berücksichtigt. Equality Saturation und E-Graphs sind dafür bekannte Werkzeuge \cite{egg}. Die erlaubten Gleichheiten müssten hier jedoch die tatsächlichen geladenen Operatorbedingungen erhalten. Die Alternative $W_6\otimes W_4$ ist beispielsweise durch den Wert $1/2$ gegen $1$ bereits ausgeschlossen.

\textbf{Prüfung.} Vergleich gegen einen guten Ward-Auswerter mit Memoisierung und dünn besetzter Arithmetik. Neue Prüfwörter enthalten gemischte Felder, Adjungierte und höhere Ordnungen. Ablationen entfernen Ladungsfilter, Symmetriekanonisierung und gemeinsame Teilrechnung getrennt.

\textbf{Möglicher Durchbruch.} Eine Größenordnung mehr exakt beantwortbare Anfragen bei gleichem Speicher und gleicher Zeit, auch außerhalb der zur Optimierung benutzten Wortfamilien. Dieses Ziel ist noch nicht erreicht. Ein begrenzter Geschwindigkeitsgewinn wäre dennoch nützlich, aber keine allgemeine neue Berechnungstheorie. Höhere Kumulanten besitzen kombinatorisch viele Partitionen; es gibt hier keine bewiesene universelle polynomielle Abkürzung.

\subsection{Baustein 2: Ein aktionsabhängiger Gedächtniscompiler}
Für bekannte endliche lineare Dynamiken sind die Bedingungen konkret. Ein Encoder $R$ erhält sämtliche Antworten, wenn für jede erlaubte Aktion $a$
\begin{equation}\label{eq:compiler}
RT_a=\widehat T_aR,\qquad\ell=\widehat\ell R.
\end{equation}
Dann folgt durch Induktion
\begin{equation}
\ell T_{a_n}\cdots T_{a_1}x
=\widehat\ell\widehat T_{a_n}\cdots\widehat T_{a_1}Rx
\end{equation}
für jedes endliche Wort. Dies ist ein wirklicher Vertrag über die Zukunft, keine Ähnlichkeit nur des momentanen Zustands.

Bei endlichen Quantenkanälen lässt sich der nötige Beobachtungsraum rückwärts schließen:
\begin{align}
\mathcal V_0&=\Span\{I,E_1,\ldots,E_m\},\\
\mathcal V_{k+1}&=\Span\!\left(\mathcal V_k\cup
\{T_a^*(E):E\in\mathcal V_k,\ a\in\mathcal A\}\right).
\end{align}
Im $d$-dimensionalen System gibt es höchstens $d^2$ unabhängige Operatorrichtungen. Für echte lineare Minimalität müssen zusätzlich die auf allen erreichbaren Zuständen unsichtbaren Richtungen entfernt werden. Der resultierende lineare Quotient muss nicht automatisch eine kleinere positive Quantendarstellung besitzen.

Für unbekannte Systeme liefert ein neues trennendes Zukunftswort die fehlende Zustandsunterscheidung. Der TFPT-Carry ist das konkrete Muster: Was heute unsichtbar ist, kann nach einer geladenen Operation wieder sichtbar werden. Die Innovation müsste in der Auswahl weniger, hochinformativer Fortsetzungen und in kontrollierter Kompression operatorwertiger oder stochastischer Prozesse liegen. Eine bloße Gegenbeispielschleife bleibt etablierte aktive Identifikation.

Auch eine angenäherte Version kann einen klaren Vertrag tragen. Sind $T_a$ und $\widehat T_a$ in den verwendeten Normen kontraktiv, gilt $\norm{RT_a-\widehat T_aR}\le\epsilon$, $\norm{\ell-\widehat\ell R}\le\delta$ und $\norm x\le1$, dann liefert die Teleskopsumme für Wortlänge $n$
\begin{equation}
\bigl|\ell T_wx-\widehat\ell\widehat T_wRx\bigr|
\le\delta+n\epsilon\norm{\widehat\ell}.
\end{equation}
Dies ist eine elementare Fehlerfortpflanzungsgrenze unter den genannten Annahmen, kein neuer allgemeiner Satz über Lernen. Sie zeigt aber, was ein praktischer Compiler versprechen kann. Eine kleine Ein-Schritt-Abweichung garantiert ohne zusätzliche Struktur keine kleine Abweichung für unbegrenzte Zukunft.

\textbf{Prüfung.} Bei deterministischen Aufgaben gegen $L^\#$; bei stochastischen Aufgaben gegen starke PSR/OOM-Verfahren und einen geeigneten Sequenzprädiktor. Verglichen werden neue längere Fortsetzungen, neue Kombinationen bekannter Aktionen, persistenter Speicher, Umweltabfragen und Wiederverwendung nach vollständigem Neustart.

\textbf{Möglicher Durchbruch.} Gleiche oder bessere Entscheidungs- und Vorhersageleistung mit wesentlich weniger dauerhaftem Speicher oder Interaktionen über mehrere verschiedenartige Aufgabenfamilien. Verschwindet der Vorteil bei gleichen Budgets und starken Baselines, ist die Überlegenheitsbehauptung zu verwerfen. Für beliebige unbeschränkte programmierbare Welten kann ein endlicher Testbestand keine universelle Zukunftsgarantie liefern.

\subsection{Baustein 3: Symmetriekompression mit erhaltenen Multiplizitäten}
Eine Symmetriedarstellung zerlegt sich als
\begin{equation}
\mathcal H=\bigoplus_\lambda(M_\lambda\otimes V_\lambda).
\end{equation}
Die Darstellung $V_\lambda$ beschreibt die Symmetriewirkung; die Multiplizitätsräume $M_\lambda$ tragen weitere tatsächlich veränderliche Information. Geladene Operationen können zwischen Sektoren wechseln. Ein korrektes Kompressionsverfahren muss deshalb sowohl die Sektoren als auch die relevanten Multiplizitäten gegenüber der ganzen Aktionsfamilie erhalten.

Symmetrische Tensor-Netze trennen solche strukturellen und freien Tensoranteile bereits \cite{singh,singhvidal}. Die hier zu entwickelnde Zusatzleistung wäre eine aktionsabhängige Auswahl der benötigten Multiplizitäten mit Antwortgarantie. Der TFPT-Befund $5\to5+70$ ist dafür eine konkrete Warnung: Eine kleine Darstellung allein ist noch kein unter späteren Operationen geschlossener Zustandsspeicher.

\textbf{Prüfung.} Vergleich gegen symmetrische Tensorverfahren und SVD-Kompression bei gleichem Speicher; eine Ablation entfernt nur die Berücksichtigung zukünftiger Aktionen. Gemessen wird der Fehler bisher ungesehener Operationsfolgen. Wenn jede Erweiterung des Handlungssatzes fast den ganzen Raum benötigt, bietet die Aufgabe wenig verlustfreie Kompression. Das ist eine Eigenschaft der Aufgabe, kein Anlass, nötige Information wegzudefinieren.

\subsection{Baustein 4: Gezielte Suche nach höheren gemeinsamen Antworten}
Der Sechserzeuge legt eine konkrete Suchregel nahe: Suche zuerst minimal ladungsneutrale Gruppen, deren echte Teilgruppen ladungsungleichgewichtig sind. Prüfe dann ihre vollständige und verbundene Antwort. So findet man Fälle, die eine rein niedrige Antwortbeschreibung übersehen kann. Minimalität der Ladungsneutralität ist eine hinreichende Auswahlhilfe für diese Zeugenklasse, keine vollständige Liste aller verbundenen Korrelationen.

Höhere Kumulanten und irreduzible gemeinsame Abhängigkeiten sind bekannte Gegenstände \cite{domino,schneidman2003}. Neu wäre ein nachweislich wirksamer Auswahlalgorithmus, der tatsächliche Operatorstruktur benutzt, um mit weniger Versuchen neue Vorhersagefehler aufzudecken. Für ein Handlungssystem könnten die Gruppen aus Kontexten und Aktionen bestehen; deren ``Ladungen'' müssten jedoch aus einer begründeten Aufgabenstruktur stammen, nicht aus einer dekorativen Physikanalogie.

\textbf{Prüfung.} Vergleich gegen zufällige Auswahl, ein paarweises Maximum-Entropie-Modell, geeignete latente Modelle und einen starken nichtlinearen Prädiktor. Paarweise Wechselwirkungen können bereits komplexe gemeinsame Statistiken erzeugen \cite{schneidman2006}. Ein nichtverschwindender höherer Kumulant beweist weder eine fundamentale Mehrkörperkraft noch Kausalität.

\textbf{Möglicher Durchbruch.} Deutlich weniger Eingriffe bis zum Auffinden unbekannter systematischer Fehler oder bis zu einer messbar besseren Entscheidung. Eine größere Sammlung berechneter Hyperkanten ohne Vorhersagewirkung wäre kein Erfolg.

\subsection{Baustein 5: Lean als Quelle ausführbarer, überprüfbarer Regeln}
Lean besitzt bereits eine Programmiersprache und Codegenerierung \cite{lean}. Die sinnvolle automatische Übernahme hängt von der Art des Originals ab:
\begin{center}
\begin{tabularx}{\linewidth}{L{0.30\linewidth}Y}
\toprule
Original in Lean & Nutzbarer Anschluss\\\midrule
Berechenbare Definition & Direkt ausführbarer Rechenkern oder Export\\
Konstruktiver Wert mit Eigenschaft & Algorithmus samt formalem Vertrag\\
Gleichheits-/Erhaltungssatz & Zulässige Optimierungsregel mit Voraussetzungen\\
Klassischer Existenzsatz & Noch kein ausführbarer Suchalgorithmus\\
Axiom oder \texttt{sorry} & Sichtbare offene Annahme beziehungsweise Beweislücke\\\bottomrule
\end{tabularx}
\end{center}

Beweise in \texttt{Prop} werden im üblichen Laufzeitcode entfernt. Eine Existenzbehauptung ohne berechenbare Konstruktion wird dadurch nicht zu einem Generator. Verifizierte Extraktion ist selbst ein etabliertes Forschungsgebiet \cite{metacoq}.

Der konkrete Weg wäre eine kleine gemeinsame Zwischensprache für rationale Arithmetik, Gitterinklusionen, Ladungen, Projektionen, Operatorwörter und Rekursion. Lean-Sätze begründen die erlaubten Umformungen. Der Quellencompiler nutzt sie nur unter ihren Bereichsbedingungen. Dadurch können wissenschaftliche Verifikationen zu wiederverwendbaren ausführbaren Regeln werden, statt in isolierten Prüfscripten zu enden.

\textbf{Prüfung.} Neue Sektoren, neue Operatoren und höhere Wortlängen müssen sowohl den formalen Vertrag als auch unabhängige Originalfälle bestehen. Ein Beweis einer falsch übersetzten Spezifikation bleibt falsch angeschlossen. Der Nutzen wäre ein schneller Kern, dessen Optimierungen zusätzliche Aufgaben abdecken und relevante Integrationsfehler verhindern; die reine Zahl grüner Lean-Sätze misst diesen Nutzen nicht.

